\section{Positioning: Where \mlir{} Ends}
\label{sec:positioning}

We frame the design space as a set of axes for \mlir{}~\cite{lattner2021mlir}. For
each axis we ask what is available today and where systems that go \emph{beyond} it
live. ``Bigger'' here does not mean a single system that dominates \mlir{} --- no such
system exists today --- but rather the research or engineering direction that pushes
that particular axis past what \mlir{} offers. Table~\ref{tab:landscape} states the
comparison.

\begin{table}[t]
  \centering
  \small
  \begin{tabularx}{\textwidth}{@{}>{\raggedright\arraybackslash}p{2.6cm} >{\raggedright\arraybackslash}X >{\raggedright\arraybackslash}X@{}}
  \toprule
  \textbf{Axis} & \textbf{\mlir{} today} & \textbf{Where ``bigger'' lives} \\
  \midrule
  Multi-level structure &
  Top form: a metalanguage; multi-level IRs are cheap &
  --- (no rival supersedes it) \\
  Declarative definitions (ops/types/attrs) &
  Top form: ODS/TableGen generates builders, accessors, verifiers &
  --- \\
  Interpreted transformations &
  Young: PDL patterns and transform-dialect scripts &
  E-graph rewriting: egg~\cite{willsey2021egg},
  egglog~\cite{zhang2023better}; tensor superoptimization:
  TASO~\cite{jia2019taso} \\
  Rewrite search and scheduling &
  Greedy, order-sensitive; fixpoints via ad-hoc pass orders &
  Equality saturation; relational e-matching~\cite{zhang2022relational};
  rule inference~\cite{nandi2021ruler} \\
  Per-compilation specialization &
  Static pipelines; passes fixed at build time; schedules left to backends &
  Autotuned schedules (TVM~\cite{chen2018tvm}, Triton~\cite{tillet2019triton});
  peephole synthesis (STOKE~\cite{schkufza2013stoke}); generative compilation
  (SPIRAL~\cite{puschel2005spiral}); staged specialization
  (Futamura~\cite{futamura1971partial}) \\
  Semantics and guarantees &
  Local per-op C++ verifiers; no formal semantics &
  Translation validation (Alive2)~\cite{lopes2021alive2}; mechanized semantics
  (Vellvm)~\cite{zhao2012vellvm}; SMT-backed contracts \\
  Incrementality &
  Whole-pass recomputation &
  Demand-driven compilation: Adapton~\cite{hammer2014adapton},
  Salsa~\cite{salsa} \\
  Runtime extensibility &
  Second-class: \texttt{ExtensibleDialect} with C++ verify callbacks;
  \texttt{.so} plugins with exact-ABI requirements &
  Data-first extension registries; sandboxed plugin hosts (WASM) \\
  Agent operability &
  Weak: extension means editing \texttt{.td}/C++ and rebuilding &
  Plan $\rightarrow$ build $\rightarrow$ load loops; machine-readable
  pass/dialect descriptions; structured diagnostics \\
  Stability and interchange &
  No core stability; dialect-level fixes (StableHLO)~\cite{stablehlo} &
  Versioned dialects with serialization discipline \\
  Schedule and hardware semantics &
  Partial: affine, linalg, CIRCT &
  Schedule-first IRs: Halide~\cite{ragankelley2013halide},
  TVM TensorIR~\cite{zheng2023tensorir}, ISL schedule trees \\
  \bottomrule
  \end{tabularx}
  \caption{Axis-by-axis comparison. ``Bigger'' names the direction that pushes an axis
  past \mlir{}, not a single competing system.}
  \label{tab:landscape}
\end{table}

Reading the table: the first two rows are MLIR's strongest ground --- the baseline for
human-authored compiler engineering --- and under the own-core decision they become
interop concerns rather than foundations. BridgeIR's claims live in the middle rows,
where it builds agent-first replacements: stable identity instead of renumbering,
runtime data-defined schemas instead of build-time ODS, transactional deltas with
provenance instead of opaque passes. The last rows we adopt as discipline (versioned
schemas, stable serialization) and per-backend work (schedule and hardware semantics).
{\color{red}\textbf{TODO:} refresh the per-axis ``BridgeIR stance'' once the design
section stabilizes.}
